\chapter{Tensor filtrations and the adelic Bost theorem}

Throughout this chapter, $S=(K,(\Omega,\mathcal A,\nu),\phi)$ is a proper
adelic curve. The field $K$ has characteristic zero. We assume either that
$\mathcal A$ is discrete or that $K$ has a common countable dense subfield
for its completed absolute values. In particular these hypotheses hold when
$K$ is countable. At Archimedean places the absolute values are the usual
real or complex ones. All bundles are finite-dimensional Hermitian adelic
vector bundles; their norms are on the completed fibres. Restrictions and
quotients always carry their actual induced norms. The notation
$\overline E\otimes\overline F$ means Hilbert tensor norms at Archimedean
places and $\varepsilon$-tensor norms elsewhere.

Write $r(E)=\dim_K E$, $\widehat{\deg}(\overline E)$ for degree, and
$\widehat\mu(\overline E)=\widehat{\deg}(\overline E)/r(E)$ for positive
rank. Semistability is defined only for positive rank and means that every
nonzero subspace has slope at most
the ambient slope; equivalently, every nonzero quotient has slope at least
the ambient slope. Write $\mathcal H_E$ for the weighted
Harder--Narasimhan filtration and $\widehat\mu_{\max}(\overline E)$ for its
largest weight when $E\ne0$. We use the Hermitian degree identities
\[
 \widehat{\deg}(\overline E)=
 \widehat{\deg}(\overline U)+\widehat{\deg}(\overline E/\overline U),
 \qquad
 \widehat{\deg}(\overline E\otimes\overline F)=
 r(F)\widehat{\deg}(\overline E)+r(E)\widehat{\deg}(\overline F),
\]
and uniqueness of the filtration with semistable grades whose slopes are
their strictly decreasing weights. These are the Hermitian slope and HN
facts recalled in Section 2 of this paper \cite{bost2026}.

\begin{definition}[Intrinsic filtration operations]
\label{def:induction-filtrations}
\source{bost-adelic-2026:page-0028}
\lean{Submodule.filtrationRightLimit, Submodule.FiltrationGraded,
Submodule.tensorFiltration, Submodule.filtrationExpectation,
Submodule.filtrationJointMultiplicity, Submodule.filtrationPairing,
Submodule.filtrationNorm}
\leanok
A finite real filtration $P$ is decreasing, exhaustive, separated,
left-continuous, and has finitely many jumps. Set
\[
 P^{>t}E=\sum_{s>t}P^sE,\qquad
 E_t=\operatorname{gr}_P^t E=P^tE/P^{>t}E,
 \qquad d(P)=\sum_t t\,r(E_t).
\]
For $E\ne0$ put $e(P)=d(P)/r(E)$. For an arbitrary family of subspaces
the same grade notation means $P^tE/(P^tE\cap P^{>t}E)$.
The filtration on the zero space is identically zero and has weighted
degree zero. The image filtration under a surjection $q$ is $P_Q^t=q(P^t)$.
The tensor filtration is
\[
 (P\otimes R)^v(E\otimes F)=\sum_{t+u\ge v}P^tE\otimes R^uF.
\]
For two finite filtrations on $E$, define their joint multiplicity by
\[
 m_{P,R}(t,u)=r(P^tE\cap R^uE)
 -r\bigl((P^{>t}E\cap R^uE)+(P^tE\cap R^{>u}E)\bigr).
\]
For positive rank set
\[
 \langle P,R\rangle=\frac1{r(E)}\sum_{t,u}tu\,m_{P,R}(t,u),
 \qquad \|P\|=\sqrt{\langle P,P\rangle}.
\]
All these sums have finite support; on the zero space their normalized
values are defined to be zero. A filtration is normalized if $e(P)=0$.
An adapted weighted basis $(b_i,x_i)$ means
$P^tE=\operatorname{span}_K\{b_i:x_i\ge t\}$ at every threshold.
\end{definition}

\begin{lemma}[A common weighted basis]
\label{lem:common-weighted-filtration-basis}
\uses{def:induction-filtrations}
\source{bost-adelic-2026:page-0028}
\lean{Submodule.exists_common_weighted_basis}
\leanok
Two finite real filtrations $P,R$ on $E$ admit a basis $(b_i)$ and
weights $(x_i,y_i)$ adapted to both. Equivalently, a vector
$\sum_i c_i b_i$ belongs to $P^tE$ precisely when $c_i\ne0$ implies
$x_i\ge t$, and belongs to $R^uE$ precisely when $c_i\ne0$ implies
$y_i\ge u$. In rank zero the basis and weight families are empty.
\end{lemma}
\begin{proof}
Write the distinct steps as increasing flags
$0=U_0\subset\cdots\subset U_h=E$ and
$0=V_0\subset\cdots\subset V_m=E$.
Suppose a basis of $U_{i-1}$ already spans every $U_{i-1}\cap V_j$
by a subset. In $U_i/U_{i-1}$, consider the increasing flag of images
of $U_i\cap V_j$. Extend a basis successively along this flag and
lift each newly added vector to $U_i\cap V_j$ at its first step.
Together with the old basis these lifts form a basis of $U_i$.
For $x\in U_i\cap V_j$, its quotient class is a combination of the
lifts chosen at steps at most $j$; subtracting that combination leaves
a vector of $U_{i-1}\cap V_j$. Thus the extended basis spans every
$U_i\cap V_j$ by the stated subsets. Induction on $i$ gives a basis
adapted to both flags. Assign to each vector its largest threshold
in each filtration; these thresholds exist by finiteness, exhaustivity,
separation and left continuity. The asserted span descriptions follow.
Uniqueness of coordinates in a basis gives the coordinate criteria.
For $E=0$ use the empty basis.
\end{proof}

\begin{lemma}[Weight formulas for expectation and pairing]
\label{lem:filtration-weight-formulas}
\uses{def:induction-filtrations,lem:common-weighted-filtration-basis}
\source{bost-adelic-2026:page-0028}
\lean{Submodule.filtrationExpectation_eq_sum_weights,
Submodule.filtrationPairing_eq_sum_weights}
\leanok
On a common weighted basis $(b_i,x_i,y_i)$ in positive rank,
\[
 e(P)=\frac{\sum_i x_i}{r(E)},\qquad
 \langle P,R\rangle=\frac{\sum_i x_i y_i}{r(E)}.
\]
These formulas compute the intrinsic quantities and are independent
of the common basis. With empty sums and normalized value zero,
the same formulas apply in rank zero.
\end{lemma}
\begin{proof}
The grade $P^tE/P^{>t}E$ has basis the classes with $x_i=t$.
Thus its dimension is the number of such indices, and summing their
weights gives the expectation formula. The intersection $P^tE\cap R^uE$
has basis the vectors with $x_i\ge t$ and $y_i\ge u$. The sum of its
two strict intersections has basis the subset for which at least one
inequality is strict. Their dimension difference counts exactly
$x_i=t,y_i=u$. Summing $tu$ times these counts gives the pairing
formula. The left sides use only the filtration subspaces, proving
independence of the basis. All dimensions and sums vanish in rank zero.
\end{proof}

\begin{lemma}[Tensor filtration weights]
\label{lem:tensor-filtration-weights}
\uses{def:induction-filtrations,lem:common-weighted-filtration-basis}
\source{bost-adelic-2026:page-0028}
\lean{Submodule.tensorFiltration_eq_span_tensorProduct}
\leanok
If $(b_i,x_i)$ and $(c_j,y_j)$ are adapted weighted bases for $P,R$,
then $(b_i\otimes c_j,x_i+y_j)$ is adapted to $P\otimes R$.
\end{lemma}
\begin{proof}
For $t+u\ge v$, expand a tensor from $P^tE\otimes R^uF$ in the
tensor basis. Every occurring pair has $x_i\ge t$, $y_j\ge u$,
hence $x_i+y_j\ge v$. This proves one span inclusion.
Conversely, if $x_i+y_j\ge v$, then
$b_i\otimes c_j\in P^{x_i}E\otimes R^{y_j}F$, one of the summands
defining $(P\otimes R)^v$. This proves the other inclusion.
\end{proof}

\begin{lemma}[Lifting the HN filtrations of the grades]
\label{lem:grade-hn-lift}
\uses{def:induction-filtrations,lem:common-weighted-filtration-basis,
lem:filtration-weight-formulas,thm:hermitian-hn-filtration,lem:degree-exact-sequence}
\source{bost-adelic-2026:page-0029,bost-adelic-2026:page-0032}
\lean{ArakelovGeometry.AdelicVectorBundle.exists_hnVariationLift}
\leanok
For every finite real filtration $P$ of a Hermitian bundle $\overline E$
there is a finite real filtration $L$ whose induced filtration on
each actual grade $E_t$ is the HN filtration of that grade with its
subquotient norms. It satisfies
\[
 e(L)=\widehat\mu(\overline E),\qquad
 r(E)\langle P,L\rangle=\sum_t t\widehat{\deg}(\overline E_t).
\]
The first formula uses the value zero in rank zero. The filtrations
$P,L$ admit a common adapted basis.
\end{lemma}
\begin{proof}
Choose an HN-adapted basis in each nonzero $E_t$, lift it to $P^tE$,
and assign its HN weights as secondary weights. Successively along the
finite flag, these lifts form a basis of $E$: a relation projects to
a relation in its last nonzero quotient, and spanning follows by
successively subtracting representatives of quotient classes.
Let $L^uE$ be the span of lifted vectors of secondary weight at least
$u$. These are finite, decreasing, exhaustive and separated cuts,
constant immediately to the left of every threshold. In the common
basis, intersection with $P^tE$ retains precisely primary weights at
least $t$; reduction modulo $P^{>t}E$ retains precisely primary weight
$t$. Thus the induced cuts on $E_t$ are its chosen HN cuts.

The sum of secondary weights over the basis of $E_t$ is
$\widehat{\deg}(\overline E_t)$ by the HN degree identity.
Summing over $t$ and using degree additivity along $P$ gives the
expectation formula by Lemma~\ref{lem:filtration-weight-formulas}.
Multiplying each grade's weight sum by its primary weight gives the
pairing formula by the same lemma. The empty construction proves
the rank-zero case.
\end{proof}

\begin{lemma}[Hermitian tensor degree]
\label{lem:induction-tensor-degree}
\uses{def:adelic-slope}
\source{bost-adelic-2026:page-0006}
\lean{ArakelovGeometry.AdelicVectorBundle.degree_orthogonalTensorPair, ArakelovGeometry.AdelicVectorBundle.slope_orthogonalTensorPair}
\leanok
For Hermitian bundles, including those of rank zero,
\[
 \widehat{\deg}(\overline E\otimes\overline F)
 =r(F)\widehat{\deg}(\overline E)+r(E)\widehat{\deg}(\overline F).
\]
For positive ranks this gives
$\widehat\mu(\overline E\otimes\overline F)
=\widehat\mu(\overline E)+\widehat\mu(\overline F)$.
\end{lemma}
\begin{proof}
For positive ranks this is the Hermitian tensor degree formula
\cite[Proposition 4.3.19, equation (4.35)]{chenmoriwaki2020}, recalled as
(8) in this paper. Its determinant norm identity identifies the norm of
the tensor determinant generator with the product of the factor
determinant norms to powers $r(F)$ and $r(E)$, respectively. Taking
negative logarithms and integrating gives the displayed equality.
Divide by $r(E\otimes F)=r(E)r(F)$ to obtain its slope form.
If one factor has rank zero, its degree and the tensor product degree
are zero, and both terms on the right vanish.
\end{proof}

\begin{lemma}[Finite contracting images of one slope]
\label{lem:finite-contracting-images}
\uses{def:adelic-slope,lem:degree-exact-sequence}
\lean{ArakelovGeometry.AdelicVectorBundle.degree_ge_of_finite_same_slope_spanning_contractions, ArakelovGeometry.AdelicVectorBundle.slope_quotient_ge_of_finite_same_slope_spanning_contractions}
\leanok
Let finitely many Hermitian bundles $\overline C_i$ be either zero or
semistable of slope $a$. Let $f_i:C_i\to B$ be contractions at every
completed place whose images span $B$. Then
\[
 \widehat{\deg}(\overline B)\ge a r(B).
\]
Every nonzero quotient of $\overline B$, with its quotient norms, has
slope at least $a$.
\end{lemma}
\begin{proof}
Order the family and set $B_j=\sum_{i\le j}f_i(C_i)$, with $B_0=0$.
The map $C_j\to B_j/B_{j-1}$ is surjective and contracting for the actual
subquotient norms. Its quotient norm is at least the target norm.
If its target is nonzero, semistability gives slope at least $a$ for
the quotient-norm bundle, and the contracting identity to the target
cannot decrease degree. Thus
$\widehat{\deg}(\overline B_j/\overline B_{j-1})\ge
a(r(B_j)-r(B_{j-1}))$. For a zero target both sides are zero.
Summing the degree identities along the flag gives the assertion.
Composing all the maps with a quotient map of $B$ preserves contraction
and spanning, so the same argument proves the quotient assertion.
\end{proof}

\begin{lemma}[Canonical tensor components]
\label{lem:tensor-grade-components}
\uses{def:induction-filtrations,lem:tensor-filtration-weights}
\source{bost-adelic-2026:page-0029}
\lean{ArakelovGeometry.NormFamily.tensorFiltrationGradedPair_pushforward_norm_le, Submodule.iSup_range_tensorFiltrationGradedQuotientPair}
\leanok
Let $P,R$ be finite filtrations on $E,F$, and let
$q:E\otimes F\twoheadrightarrow Q$ have the original tensor quotient norms.
For $t+u=\gamma$ the canonical map
\[
 E_t\otimes F_u\longrightarrow Q_\gamma,
 \qquad Q_\gamma=\operatorname{gr}_{(P\otimes R)_Q}^{\gamma}Q,
\]
is contracting at every completed place. The images of these maps span
$Q_\gamma$. This includes $q=\operatorname{id}$.
\end{lemma}
\begin{proof}
The map sends the tensor of classes of $x\in P^tE$ and $y\in R^uF$
to the class of $q(x\otimes y)$. Changing $x$ or $y$ by a vector of
strictly higher weight changes its image by a vector of weight strictly
greater than $\gamma$, so the map is well defined.
At an Archimedean place, the quotient maps defining $E_t,F_u$ have
orthogonal isometric sections. Their Hilbert tensor product is isometric;
the subsequent restriction and quotient maps are contracting.
At a non-Archimedean place, for every $c>1$ the finite-dimensional quotient
maps admit linear sections of operator norm at most $c$, by choosing
arbitrarily nearly orthogonal quotient bases and their nearly minimal lifts.
The $\varepsilon$-tensor map of these sections has norm at most $c^2$.
The same subsequent maps are contracting. Since this holds for every
$c>1$, the canonical component has norm at most one. Zero grades give
the zero map and require no section.
Finally choose adapted bases for $P,R$. Tensor basis vectors of weight
at least $\gamma$ span $(P\otimes R)^\gamma$; those of weight greater
than $\gamma$ vanish in the grade. The remaining classes, of weight
exactly $\gamma$, are images of the displayed component maps and span
the quotient grade.
\end{proof}

\begin{lemma}[Tensor HN identification]
\label{lem:tensor-hn-identification}
\uses{lem:finite-contracting-images,lem:tensor-grade-components,thm:hermitian-hn-filtration,lem:hn-uniqueness,lem:degree-exact-sequence,lem:induction-tensor-degree,lem:tensor-filtration-weights,lem:filtration-weight-formulas}
\source{bost-adelic-2026:page-0029}
\lean{ArakelovGeometry.AdelicVectorBundle.hermitianHNFiltration_orthogonalTensorPair_of_semistable_grades}
\leanok
Suppose that the tensor product of every pair of nonzero HN grades of
$\overline E,\overline F$ is semistable. Then
\[
 \mathcal H_{E\otimes F}=\mathcal H_E\otimes\mathcal H_F.
\]
Zero-rank factors are allowed, with the zero-filtration convention.
\end{lemma}
\begin{proof}
Put $P=\mathcal H_E\otimes\mathcal H_F$. Every nonzero source
$E_t\otimes F_u$ of a component of weight $\gamma=t+u$ is semistable
of slope $\gamma$, by the hypothesis and the tensor degree formula.
The two preceding lemmas give
\[
 \widehat{\deg}(\operatorname{gr}_P^\gamma(\overline E\otimes\overline F))
 \ge \gamma\,r(\operatorname{gr}_P^\gamma(E\otimes F)),
\]
and the corresponding lower bound $\gamma$ for every nonzero quotient
of that actual grade. Choose HN-adapted bases with weights $x_i,y_j$.
The sum of the right sides is
\[
 \sum_{i,j}(x_i+y_j)=r(F)\widehat{\deg}(\overline E)
  +r(E)\widehat{\deg}(\overline F)
  =\widehat{\deg}(\overline E\otimes\overline F).
\]
The sum of the left sides is the same by degree additivity. These are
finite sums of inequalities, so each is an equality. Every nonzero grade
therefore has slope $\gamma$, and its quotient bounds prove semistability.
The weights strictly decrease along its nonzero grades; HN uniqueness
gives the identity. If a factor is zero, both filtrations are identically zero.
This proof uses only the individual component contractions, rather than
the contraction of the whole orthogonal sum asserted in (44) of this paper.
\end{proof}

\begin{lemma}[HN images and degree]
\label{lem:finite-hn-image-degree}
\uses{lem:finite-contracting-images,def:induction-filtrations,thm:hermitian-hn-filtration,lem:degree-exact-sequence}
\source{bost-adelic-2026:page-0029,bost-adelic-2026:page-0030}
\lean{ArakelovGeometry.AdelicVectorBundle.hermitianHN_imageFiltration_finite, ArakelovGeometry.AdelicVectorBundle.finsum_hermitianHN_imageFiltration_le_degree}
\leanok
Let $f_i:C_i\to B$ be a finite family of contractions with spanning
images. Define $H^sB=\sum_i f_i(\mathcal H_{C_i}^s C_i)$. Then
$H$ is a finite filtration and $d(H)\le\widehat{\deg}(\overline B)$.
\end{lemma}
\begin{proof}
Exhaustivity, separation and left continuity follow by taking common
bounds for the finitely many source filtrations. Moreover
$H^{>s}=\sum_i f_i(\mathcal H_{C_i}^{>s})$, by interchanging the sums.
Consequently the canonical maps from the HN grades at $s$ jointly span
$\operatorname{gr}_H^s B$ and remain contracting for the actual
subquotient norms. Each nonzero source grade is semistable of slope $s$.
The finite-image lemma bounds the degree of this grade below by its
rank times $s$. Sum these inequalities over the finitely many nonzero
grades and use degree additivity.
\end{proof}

\begin{lemma}[One-sided HN variation]
\label{lem:component-hn-variation}
\uses{lem:finite-hn-image-degree,lem:tensor-grade-components,def:induction-filtrations,lem:grade-hn-lift,lem:tensor-filtration-weights,lem:filtration-weight-formulas}
\source{bost-adelic-2026:page-0030,bost-adelic-2026:page-0031}
\lean{ArakelovGeometry.AdelicVectorBundle.tensorFiltrationQuotient_eventual_hn_upper}
\leanok
Let $q:E\otimes F\twoheadrightarrow Q\ne0$, let $P,R$ be finite
filtrations, and choose lifts of their grade HN filtrations with common
adapted bases. Suppose all pairs of grades have rank sum strictly less
than $r(E)+r(F)$ and the tensor HN identity holds for every such pair.
Let $P(\varepsilon),R(\varepsilon)$ replace primary weights $x$ by
$x+\varepsilon y$, where $y$ is the chosen secondary weight. In each
primary quotient grade $Q_\gamma$, let $H_\gamma$ be the sum of the
images of the tensor HN filtrations of $E_t\otimes F_u$, $t+u=\gamma$.
Put $r=r(Q)$ and $\delta=r^{-1}\sum_\gamma d(H_\gamma)$. Then
\[
 \delta\le\widehat\mu(\overline Q),\qquad
 e((P(\varepsilon)\otimes R(\varepsilon))_Q)
 \le e((P\otimes R)_Q)+\varepsilon\delta
\]
for all sufficiently small $\varepsilon>0$.
\end{lemma}
\begin{proof}
The tensor HN hypothesis identifies each source secondary filtration
with its actual HN filtration. Apply the preceding lemma to the canonical
components in each $Q_\gamma$. It gives
$d(H_\gamma)\le\widehat{\deg}(\overline Q_\gamma)$.
Summing along the primary quotient flag and dividing by $r$ proves
the first inequality.

For the second inequality, write $z$ for the tensor basis vectors and
$p_z,s_z$ for their primary and secondary weights. Let $\mathcal B$ be
the finite nonempty family of $r$-element subsets whose images form a
basis of $Q$. For arbitrary weights $w_z$, the weighted degree of the
image filtration is $\max_{I\in\mathcal B}\sum_{z\in I}w_z$.
Indeed, list the distinct weights $c_1>\cdots>c_h$, and let $d_i$ be
the dimension of the span of images of weights at least $c_i$.
For any image basis, its number $n_i$ of such vectors satisfies $n_i\le d_i$,
so its total weight is at most
$c_h r+\sum_{i<h}(c_i-c_{i+1})d_i$.
Successively extending a basis along these spans attains this bound,
which is exactly the filtration degree.

Put $\alpha_I=\sum_{z\in I}p_z$, $\beta_I=\sum_{z\in I}s_z$,
$M=\max_I\alpha_I$, and $\mathcal B_0=\{I:\alpha_I=M\}$.
Finiteness implies that, for sufficiently small $\varepsilon>0$,
only $I\in\mathcal B_0$ maximize $\alpha_I+\varepsilon\beta_I$:
if there are other sets, take the minimum positive primary gap $\Delta$
and choose $\varepsilon<\Delta/(2\max_I|\beta_I|+1)$.
Every $I\in\mathcal B_0$ attains all the preceding primary rank bounds.
Its vectors of each primary weight $\gamma$ therefore give a basis of
$Q_\gamma$. The class of a tensor vector of secondary weight $s_z$
belongs to $H_\gamma^{s_z}$, by the definition of the lifts and the
component images. In any filtered space, the sum of assigned weights
of a basis whose vectors lie in their assigned cuts is at most the
filtration degree: at each threshold their number is at most the
dimension of the cut. Below all assigned and filtration weights these
counts both equal the full dimension, so the same telescoping formula
proves the bound, including negative weights.
Applying it in every primary grade gives
$\beta_I\le\sum_\gamma d(H_\gamma)=r\delta$.
The maximum formula now gives the desired inequality after division by $r$.
This is the one-sided form of Lemma 7.2 needed below; it does not assert
the stronger affine identity or a contraction from an orthogonal direct sum.
\end{proof}

\begin{lemma}[Original-field geometric optimizer]
\label{lem:induction-geometric-optimizer}
\uses{def:induction-filtrations,lem:filtration-weight-formulas}
\source{bost-adelic-2026:page-0031}
\lean{Submodule.exists_normalized_tensor_quotient_optimizer_of_not_geometricallySemistable}
\leanok
Let $E,F,Q$ have positive ranks and let
$q:E\otimes F\twoheadrightarrow Q$ be geometrically unstable. Here
geometric semistability means semistability of the dual Pl\"ucker point
of $Q^\vee\subset(E\otimes F)^\vee$ for the product of the special
linear groups after passage to an algebraic closure. There exist
normalized finite filtrations $P,R$ over $K$ such that
\[
 \|P\|^2+\|R\|^2=1,\qquad e((P\otimes R)_Q)<0,
\]
and, for every normalized finite pair $A,B$ not both trivial,
\[
 e((P\otimes R)_Q)\le
 \frac{e((A\otimes B)_Q)}{\sqrt{\|A\|^2+\|B\|^2}}.
\]
\end{lemma}
\begin{proof}
Apply the original-field geometric instability theorem in its binary
case \cite[Theorem 5.4.3]{chenmoriwaki2020}, as used in Section 7.3 of
this paper. It gives a normalized nontrivial pair minimizing the
displayed ratio, with negative minimum. Its norm $L$ is positive.
Divide both sets of weights by $L$. Expectations scale by $1/L$ and
the sum of squared norms scales by $1/L^2$. The new pair is normalized,
has squared norm one, and has quotient expectation equal to the same
negative minimum. The minimizing property gives exactly the stated
comparison with every $A,B$. The original-field conclusion of the
instability theorem is used here; no extension of the adelic norms to
an algebraic closure is made.
\end{proof}

\begin{lemma}[Unstable quotient bound]
\label{lem:unstable-quotient-bound}
\uses{lem:component-hn-variation,def:induction-filtrations,lem:degree-exact-sequence,lem:induction-geometric-optimizer,lem:grade-hn-lift,lem:filtration-weight-formulas}
\source{bost-adelic-2026:page-0031,bost-adelic-2026:page-0032}
\lean{ArakelovGeometry.AdelicVectorBundle.slope_sum_le_tensor_quotient_of_not_geometricallySemistable}
\leanok
Let $\overline E,\overline F$ be positive-rank semistable Hermitian
bundles and $q:E\otimes F\twoheadrightarrow Q\ne0$ a quotient with
its original tensor quotient norms. Assume that for every pair of
real-indexed families $P,R$ of subspaces and every $t,u$ with
$r(E_t)+r(F_u)<r(E)+r(F)$, the actual tensor HN identity holds on
$E_t\otimes F_u$. If the quotient is not geometrically semistable,
then
\[
 \widehat\mu(\overline Q)\ge
 \widehat\mu(\overline E)+\widehat\mu(\overline F).
\]
Here geometric semistability means semistability of the dual Pl\"ucker
point of $Q^\vee\subset(E\otimes F)^\vee$ for the product of the
special linear groups, after passage to an algebraic closure.
\end{lemma}
\begin{proof}
The geometric instability theorem supplies an original-field pair of
normalized finite filtrations $P,R$, not both trivial, minimizing
\[
 \Theta(P,R)=\frac{e((P\otimes R)_Q)}
 {\sqrt{\|P\|^2+\|R\|^2}}.
\]
Write $L=\sqrt{\|P\|^2+\|R\|^2}>0$ and $c=\Theta(P,R)<0$.
At least one normalized filtration has more than one nonzero grade;
otherwise both have their sole weight zero. Each of its grades has
rank strictly less than the whole factor. The other factor's grades
have rank at most that factor's rank. Thus every grade pair has
strictly smaller rank sum, including pairs with a zero grade.

Choose grade HN lifts $L_E,L_F$. Degree additivity gives
$e(L_E)=\widehat\mu(\overline E)$ and
$e(L_F)=\widehat\mu(\overline F)$. Set
$\beta=\langle P,L_E\rangle+\langle R,L_F\rangle$.
For decreasing primary weights $t_1>\cdots>t_h$, put
$U_i=P^{t_i}E$, $U_0=0$, and $m=r(E)$. Then
\[
\begin{split}
 m\langle P,L_E\rangle
 &=\sum_i t_i\widehat{\deg}(\overline U_i/\overline U_{i-1})\\
 &=t_h\widehat{\deg}(\overline E)
   +\sum_{i<h}(t_i-t_{i+1})\widehat{\deg}(\overline U_i)\\
 &\le\widehat\mu(\overline E)
   \left(t_hm+\sum_{i<h}(t_i-t_{i+1})r(U_i)\right)=0.
\end{split}
\]
The inequality is semistability; the last equality is normalization.
The identical calculation for $F$ gives $\beta\le0$.

Subtract the two mean slopes from the corresponding lift weights and
perturb in these centered directions. The pair remains normalized.
The preceding variation lemma gives a number
$\delta\le\widehat\mu(\overline Q)$ and, writing
$a=\widehat\mu(\overline E)+\widehat\mu(\overline F)$, an upper bound
$e((P\otimes R)_Q)+\varepsilon(\delta-a)$ for the perturbed numerator.
The perturbed denominator is
$L+\varepsilon\beta/L+o(\varepsilon)$, since centering the lifts does
not change their inner products with normalized $P,R$.
It stays positive for small $\varepsilon$. Minimality yields
\[
 c\bigl(L+\varepsilon\beta/L+o(\varepsilon)\bigr)
 \le cL+\varepsilon(\delta-a).
\]
Subtract $cL$, divide by $\varepsilon>0$, and take the limit. Thus
$\delta-a\ge c\beta/L\ge0$, which proves the assertion.
\end{proof}

\begin{proposition}[Assembly from the geometric quotient inequality]
\label{prop:conditional-tensor-hn-assembly}
\uses{lem:unstable-quotient-bound,lem:tensor-hn-identification,thm:hermitian-hn-filtration,lem:induction-tensor-degree}
\source{bost-adelic-2026:page-0031,bost-adelic-2026:page-0032}
\lean{ArakelovGeometry.AdelicVectorBundle.tensorHN_and_semistable_of_geometric_quotient_bound}
\leanok
Assume the following inequality: for every pair of positive-rank
semistable Hermitian bundles and every nonzero geometrically semistable
quotient of their tensor product, with the original quotient norms,
\[
 \widehat\mu(\overline Q)\ge
 \widehat\mu(\overline E)+\widehat\mu(\overline F).
\]
Then tensor HN compatibility holds for all pairs of Hermitian bundles,
including zero-rank bundles, and the tensor product of two semistable
bundles is semistable.
\end{proposition}
\begin{proof}
Use strong induction on a bound $n$ for the sum of the two ranks,
proving both assertions simultaneously for all pairs of rank sum at most $n$.
First fix a semistable pair of such rank sum. For each nonzero quotient,
the geometrically semistable case follows from the displayed assumption.
In the other case use the unstable quotient lemma. Its hypothesis holds
for every pair of arbitrary subspace families: whenever the grade ranks
have sum smaller than the fixed pair's rank sum, they have sum smaller
than $n$, so the HN identity is the first conclusion of the induction
hypothesis for those actual subquotient bundles. The quotient inequality
and the tensor slope formula therefore prove semistability.

Next take any pair with rank sum at most $n$. Each nonzero HN grade is
semistable, and a pair of grades has rank sum at most $n$. The semistable
assertion just proved, which was uniform over all pairs of this bound,
applies even when both grades have full rank. Tensor HN identification
now proves the other assertion. If a factor has rank zero, every tensor
grade is zero and the identity holds by the zero-filtration convention.
For $n=0,1$ there are no pairs of positive-rank semistable factors, so
this also proves the initial induction steps. No smaller-rank hypothesis
is retained in the conclusion.
\end{proof}

\begin{theorem}[Tensor HN compatibility]
\label{thm:tensor-hn-compatibility}
\uses{prop:conditional-tensor-hn-assembly,thm:geometric-quotient-bound}
\source{bost-adelic-2026:page-0032}
\notready
Suppose in addition that $K$ is countable. For Hermitian adelic bundles on $S$,
$\mathcal H_{E\otimes F}=\mathcal H_E\otimes\mathcal H_F$.
The tensor product of positive-rank semistable bundles is semistable.
\end{theorem}
\begin{proof}
The geometric quotient inequality supplies exactly the assumption of
the preceding proposition, for the original quotient norms. Apply that
proposition to the given bundles.
\end{proof}

\begin{lemma}[Maximal slope from the actual tensor filtration]
\label{lem:maximal-slope-from-tensor-hn}
\uses{def:induction-filtrations,thm:hermitian-hn-filtration,lem:induction-tensor-degree,lem:tensor-filtration-weights}
\source{bost-adelic-2026:page-0032}
\lean{ArakelovGeometry.AdelicVectorBundle.maximalSlope_orthogonalTensorPair_of_hnFiltration_eq}
\leanok
If $E,F\ne0$ and $\mathcal H_{E\otimes F}=\mathcal H_E\otimes\mathcal H_F$,
then
\[
 \widehat\mu_{\max}(\overline E\otimes\overline F)=
 \widehat\mu_{\max}(\overline E)+\widehat\mu_{\max}(\overline F).
\]
\end{lemma}
\begin{proof}
Let $a,b$ be the largest HN weights of the two factors. Every weight of
an adapted basis is at most the corresponding largest weight. Thus the
tensor filtration has zero cut at every $v>a+b$. By the asserted HN
identity, its largest HN slope is at most $a+b$.
For the reverse inequality, the first HN subspaces $E_1,F_1$ have
slopes $a,b$. Their tensor product has slope $a+b$, and its natural
inclusion into $E\otimes F$ is contracting for the stated tensor norms.
Its image, with induced norms, has degree at least that of
$\overline E_1\otimes\overline F_1$. Hence the maximal slope of the
ambient tensor product is at least $a+b$. The ranks of $E_1,F_1$ are
positive, so this supplies a nonzero subspace as required.
\end{proof}

\begin{lemma}[The one-factor tensor isometry]
\label{lem:finite-tensor-singleton-isometry}
\uses{def:adelic-slope}
\lean{ArakelovGeometry.NormFamily.orthogonalTensor_norm_subsingletonEquiv,
ArakelovGeometry.AdelicVectorBundle.slopes_orthogonalTensor_subsingleton}
\leanok
For a family with one index $i_0$, the canonical identification
$\bigotimes_i E_i\simeq E_{i_0}$ preserves every completed norm.
It preserves the mean slope and the maximal slope.
\end{lemma}
\begin{proof}
Every one-factor completed tensor is a pure tensor, and its tensor norm
equals the norm of its sole factor. Thus the canonical linear equivalence
is an isometry. Its determinant map preserves determinant norms, hence
degrees and ranks. It also gives a bijection of subspaces preserving
their restricted norms, hence the supremum of their slopes. This proves
both slope identities.
\end{proof}

\begin{lemma}[Finite tensor reindexing]
\label{lem:finite-tensor-reindex-isometry}
\uses{def:adelic-slope}
\lean{ArakelovGeometry.NormFamily.orthogonalTensor_norm_reindex,
ArakelovGeometry.AdelicVectorBundle.slopes_orthogonalTensor_reindex}
\leanok
For a bijection $\sigma:I\simeq J$ of finite index sets, the canonical map
\[
 \bigotimes_{i\in I} E_i\longrightarrow
 \bigotimes_{j\in J}E_{\sigma^{-1}(j)}
\]
is an isometry at every completed place. It preserves mean and maximal
slopes.
\end{lemma}
\begin{proof}
At an Archimedean place choose orthonormal bases in the factors.
The tensor norm is the square root of the sum of squared absolute values
of the tensor-basis coefficients. Reindexing permutes those coefficients,
so it preserves this sum. At a non-Archimedean place the injective norm
is the supremum of the absolute values of product-dual evaluations divided
by the products of the dual norms, with zero denominators contributing
zero. Reindexing bijects these families of functionals and preserves both
the evaluation and the denominator, so it preserves the supremum. In both cases
the equalities hold for every completed vector. Determinants and
restricted subspaces are preserved isometrically, giving the two slope
identities as in the one-factor case.
\end{proof}

\begin{lemma}[Finite tensor grouping]
\label{lem:finite-tensor-grouping-isometry}
\uses{def:adelic-slope}
\lean{ArakelovGeometry.NormFamily.orthogonalTensor_norm_tmulEquivDep,
ArakelovGeometry.AdelicVectorBundle.slopes_orthogonalTensor_sum}
\leanok
For finite disjoint index sets $I,J$ and a family on $I\sqcup J$,
the canonical grouping map
\[
 \left(\bigotimes_{i\in I}E_i\right)\otimes
 \left(\bigotimes_{j\in J}E_j\right)
 \longrightarrow\bigotimes_{k\in I\sqcup J}E_k
\]
is an isometry at every completed place. It preserves mean and maximal
slopes.
\end{lemma}
\begin{proof}
At an Archimedean place the tensor products of factor orthonormal bases
are orthonormal in both presentations. Grouping matches their vectors
and their coefficients, so preserves every squared-norm sum.

At a non-Archimedean place fix $0<\alpha<1$ and choose an
$\alpha$-orthogonal basis in each nonzero factor. Write $N=|I|+|J|$ and,
for a tensor with coefficients $z_{\mathbf i}$, put
\[
 M(z)=\max_{\mathbf i}|z_{\mathbf i}|
                    \prod_k\|b_{k,i_k}\|.
\]
For the ungrouped injective norm, the ultrametric inequality and the
crossnorm formula give the upper bound $M(z)$. The coordinate functional
dual to $b_{k,i_k}$ has norm at most
$\alpha^{-1}\|b_{k,i_k}\|^{-1}$; evaluating their product gives the
lower bound $\alpha^N M(z)$. The same argument on each block shows that
its tensor basis is $\alpha^{|I|}$- or $\alpha^{|J|}$-orthogonal,
with basis-vector norms the corresponding products. Applying the two-factor
argument to those block bases bounds the grouped norm between
$\alpha^N M(z)$ and $M(z)$ as well. Consequently the two norms differ
by a ratio between $\alpha^N$ and $\alpha^{-N}$. This comparison holds
for every $\alpha<1$, although its bases may vary, so passage to
$\alpha\uparrow1$ proves equality. If a factor is zero both tensor
spaces vanish. Empty blocks use the scalar field with its absolute value,
for which the same calculation has an empty product equal to one.

The determinant and subspace argument again proves preservation of
both slopes. Successive grouping and reindexing therefore account for
any enumeration or parenthesization of a finite tensor family.
\end{proof}

\begin{lemma}[Mean slope of a finite orthogonal tensor]
\label{lem:finite-tensor-mean-slope}
\uses{lem:finite-tensor-singleton-isometry,lem:finite-tensor-reindex-isometry,lem:finite-tensor-grouping-isometry,lem:induction-tensor-degree,def:adelic-slope}
\source{bost-adelic-2026:page-0006}
\lean{ArakelovGeometry.AdelicVectorBundle.finrank_piTensorProduct_eq_prod, ArakelovGeometry.AdelicVectorBundle.finrank_piTensorProduct_pos, ArakelovGeometry.AdelicVectorBundle.slope_orthogonalTensorOfMeasurableBase}
\leanok
For every finite nonempty set $I$ and family
$(\overline E_i)_{i\in I}$ of positive-rank Hermitian bundles, the actual
finite orthogonal tensor bundle $\overline T=\bigotimes_{i\in I}\overline E_i$
satisfies
\[
 r(T)=\prod_{i\in I}r(E_i)>0,\qquad
 \widehat\mu(\overline T)=\sum_{i\in I}\widehat\mu(\overline E_i).
\]
\end{lemma}
\begin{proof}
We iterate the binary degree formula (8) and its slope form (9) in this
paper \cite{bost2026}, identifying the completed tensor norms at each step.
The tensor products of basis vectors form a basis indexed by the product
of the factor basis sets, proving the rank formula and positivity.
For a singleton, the canonical tensor map to its unique factor is an
isometry: every completed tensor is a one-factor pure tensor and its norm
is that factor's norm. Thus the slope formula holds for one factor.

For a larger family, choose $i_0\in I$ and put $J=I\setminus\{i_0\}$,
which is nonempty. The canonical reindexing and grouping map identifies
the finite tensor with
$\overline E_{i_0}\otimes(\bigotimes_{j\in J}\overline E_j)$ isometrically
at every completed place by
Lemmas~\ref{lem:finite-tensor-reindex-isometry} and
\ref{lem:finite-tensor-grouping-isometry}.
This is an identification of the full completed norms on all vectors.
An isometry preserves ranks and determinant norms, hence degrees and mean
slopes. The two factors in this grouping have positive rank. The binary
mean-slope formula, followed by induction on $|I|$, gives the displayed sum.
\end{proof}

\begin{lemma}[Finite maximal slopes from binary additivity]
\label{lem:finite-tensor-slopes-of-binary}
\uses{lem:finite-tensor-singleton-isometry,lem:finite-tensor-reindex-isometry,lem:finite-tensor-grouping-isometry,lem:finite-tensor-mean-slope,def:adelic-slope}
\source{bost-adelic-2026:page-0001,bost-adelic-2026:page-0032}
\lean{ArakelovGeometry.AdelicVectorBundle.maximalSlope_orthogonalTensorOfMeasurableBase_of_binary_additivity}
\leanok
Assume that maximal slopes are additive for every pair of positive-rank
Hermitian bundles on $S$. Then every finite nonempty family as above satisfies
\[
 \widehat\mu_{\max}(\overline T)=
 \sum_{i\in I}\widehat\mu_{\max}(\overline E_i).
\]
\end{lemma}
\begin{proof}
Induct on the size of $I$. For a singleton the canonical identification
with its factor is an isometry. For $|I|>1$, choose $i_0$ and set
$J=I\setminus\{i_0\}$. Reindexing and grouping identify the finite tensor
isometrically with $\overline E_{i_0}\otimes\bigotimes_{j\in J}\overline E_j$.
The preceding lemma ensures that the latter factor has positive rank.
An isometry bijects nonzero subspaces with isometric restricted bundles,
so it preserves their slope supremum. Apply the binary hypothesis to the
two factors and then the induction hypothesis to $J$. This gives the sum
over $I$.
\end{proof}

\begin{corollary}[Finite tensor semistability from binary additivity]
\label{cor:finite-tensor-semistability-of-binary}
\uses{lem:finite-tensor-slopes-of-binary,lem:finite-tensor-mean-slope,def:adelic-slope}
\source{bost-adelic-2026:page-0001,bost-adelic-2026:page-0032}
\lean{ArakelovGeometry.AdelicVectorBundle.isMeanSlopeSemistable_orthogonalTensorOfMeasurableBase_of_binary_additivity}
\leanok
Under the binary maximal-slope hypothesis of the preceding lemma, the
orthogonal tensor product of a finite nonempty family of positive-rank
semistable Hermitian bundles is semistable.
\end{corollary}
\begin{proof}
For each semistable factor, the slope of every nonzero subspace is at most
its mean slope, and the whole space is an allowed subspace. Thus its maximal
slope equals its mean slope. The preceding two lemmas give
$\widehat\mu_{\max}(\overline T)=\widehat\mu(\overline T)$.
Every nonzero subspace of $T$ has slope at most its maximal slope and hence
at most its mean slope. Since $T$ has positive rank, this is semistability.
\end{proof}

\begin{theorem}[Adelic Bost theorem]
\label{thm:bost-adelic}
\uses{thm:tensor-hn-compatibility,lem:maximal-slope-from-tensor-hn,lem:finite-tensor-slopes-of-binary,cor:finite-tensor-semistability-of-binary}
\source{bost-adelic-2026:page-0001,bost-adelic-2026:page-0032}
\notready
Let $K$ be countable of characteristic zero and $S$ proper, with the
usual Archimedean absolute values. For every pair of positive-rank
Hermitian adelic vector bundles, with Hilbert/\(\varepsilon\) tensor norms,
\[
 \widehat\mu_{\max}(\overline E\otimes\overline F)=
 \widehat\mu_{\max}(\overline E)+\widehat\mu_{\max}(\overline F).
\]
For every finite nonempty family of positive-rank factors the maximal
slope of its tensor product is the sum of their maximal slopes. In
particular a tensor product of semistable factors is semistable.
\end{theorem}
\begin{proof}
Countability supplies the common countable dense subfield hypothesis
by taking $K$ itself. Apply tensor HN compatibility, then the preceding
maximal-slope lemma. This supplies the binary hypothesis of
Lemma~\ref{lem:finite-tensor-slopes-of-binary}, which proves the finite-family
identity on the actual tensor norms. Apply
Corollary~\ref{cor:finite-tensor-semistability-of-binary} for semistable factors.
\end{proof}
